\section{Reproducibility notes}
\label{sec:repro}

\paragraph{Training details.} We set free bits to 0 rather than DreamerV3's reference default of 1.0.
The measured KL divergence settles near 1.8 nats, well above collapse, so the floor is inactive in
this setting. We report the change because otherwise describing the setup as the published defaults
would be inaccurate. All three trained models pass the acceptance checks; none were discarded.

\paragraph{Extraction data.} Training uses trajectories 0--203 from the 256-trajectory dataset.
Extraction uses the remaining 52 trajectories. We fit $C$ on those 52 held-out trajectories and
evaluate the Section~\ref{sec:intervention} intervention on the same 52. The recovered and
random-constraint arms use the same data split.

\paragraph{Calibrating the displacement.} The sensitivity analysis in Section~\ref{sec:leverage}
needs a perturbation large enough to produce a measurable nonlinear effect. At $\epsilon = 0.05|z|$,
damage is about $-0.5\%$ of baseline and its sign varies across directions. Because imagination is
deterministic, repeated unperturbed rollouts are identical, so these small values do not come from
sampling noise.

For the leading direction, damage changes by $-0.5\%$, $+3.1\%$, $+18.1\%$, $+46.4\%$, and
$+128.8\%$ at $\epsilon/|z| = 0.05$, $0.10$, $0.25$, $0.50$, and $1.00$. We use
$\epsilon = 0.25|z|$, the smallest perturbation whose damage exceeds 10\% of baseline while the
response still grows smoothly.

\paragraph{Scoring the intervention over a grid.} We score the slope of rollout error over the fixed
$\alpha$ grid rather than the minimum value on that grid. Taking the minimum would select the most
favorable of five points after observing the curve and can misrepresent the overall response. In one
diagnostic run, that rule would have reported a 5.9\% improvement even though error rose
monotonically to twice baseline across the sweep. The slope avoids that failure mode.

\paragraph{Testing frequency weighting.} The weighted and unweighted integral predictors are nested,
so comparing their standalone $R^2$ values is not enough. A difference of $0.019$ in $R^2$, for
example, does not establish which predictor is better unless we compare that increment with its
sampling uncertainty. The weighted predictor leads in 6 of 9 seed--horizon combinations, but the
nested bootstrap test in Section~\ref{sec:ablations} finds a significant incremental contribution in
only 1 of 9. We therefore report the frequency weighting as unresolved rather than supported or
refuted.

\paragraph{Extraction dimension and residual decomposition.} Figure~\ref{fig:ldsweep} reports
recovery and pairing residual across extraction dimension. Figure~\ref{fig:lowvar} shows where energy
information, latent flow, and pairing residual lie across the high- and lower-variance principal
directions.

\begin{figure}[h]
\centering
\begin{minipage}[t]{0.36\textwidth}
\centering
\includegraphics[width=\textwidth]{figures/fig2_ld_sweep.pdf}
\caption{Energy recovery and pairing residual versus extraction dimension. Lines show medians over
the three models and points show individual models. Both quantities are dimensionless and plotted on
the same 0--1 axis.}
\label{fig:ldsweep}
\end{minipage}\hfill
\begin{minipage}[t]{0.58\textwidth}
\centering
\includegraphics[width=\textwidth]{figures/fig3_low_variance.pdf}
\caption{\textbf{(a)} Energy-probe $R^2$ for the top six principal directions versus directions
7--12, shown separately for each model. \textbf{(b)} Fraction of latent flow and pairing residual
carried by each group, using medians across the three models.}
\label{fig:lowvar}
\end{minipage}
\end{figure}
